% !TeX program = xelatex
% !BIB program = biber
\documentclass[light]{lutbeamer}
\usepackage{pgfpages}
\usepackage{ragged2e}
\usepackage{booktabs}
\usepackage{amssymb}
\usepackage{fontawesome5}
\usepackage{tikz}
\usepackage{pgfplots}
\pgfplotsset{compat=1.18}
\addbibresource{references.bib}

% ===================================================================
% METADATA (DERS KIMLIK BILGILERI)
% ===================================================================
\setdepartment{OSTIM Teknik Üniversitesi}
\institute{}
\author[Mühendislik Fakültesi]{\\
Prof.~Dr.~Yalçın ATA \\
Prof.~Dr.~Arif DEMİR \\
Dr.~Öğr.~Üyesi~Muhammed ELMNEFİ \\
Dr.~Öğr.~Üyesi~Haydar KILIÇ \\
Dr.~Öğr.~Üyesi~Emel GÜVEN \\
Dr.~Öğr.~Üyesi~Ufuk ASIL \\
Öğr.~Gör.~Sema ÇİFTÇİ
}
\title{MATH 204 -- Mühendisler için Olasılık ve İstatistik}
\subtitle{Hafta X: Beklenen Değer, Varyans ve Standart Sapma}
\date{2025--2026 Bahar Dönemi}

% ===================================================================
% AUTOMATIC TABLE OF CONTENTS (PRESENTATION FLOW)
% ===================================================================
\setcounter{tocdepth}{2}

\AtBeginSection[]{
  \begin{frame}{Sunum Planı}
    \scriptsize
    \begin{columns}[T]
      \begin{column}{0.48\textwidth}
        \tableofcontents[sections={1-4}, currentsection]
      \end{column}
      \begin{column}{0.48\textwidth}
        \tableofcontents[sections={5-10}, currentsection]
      \end{column}
    \end{columns}
  \end{frame}
}

\begin{document}

\setbeamertemplate{navigation symbols}{}
\setbeamertemplate{footline}{
  \begin{beamercolorbox}[wd=\paperwidth, ht=10pt, dp=1pt]{footline}
    \usebeamerfont{footline}
    \begin{tikzpicture}[remember picture,overlay]
      \node[anchor=south west, xshift=10mm, yshift=.4mm]
      at (current page.south west)
      {\insertshortauthor,~\department};
    \end{tikzpicture}
    \begin{tikzpicture}[remember picture,overlay]
      \node[anchor=south east, xshift=-5mm, yshift=0.5mm]
      at (current page.south east)
      {\insertframenumber \,/\, \inserttotalframenumber};
    \end{tikzpicture}
  \end{beamercolorbox}
}

% ===================================================================
% KAPAK SLAYTLARI
% ===================================================================

% Gorsel Tam Kapak (Logosuz)
\begin{frame}<article:0>[plain,noframenumbering]
  \begin{tikzpicture}[remember picture,overlay]
    \node[anchor=center, at=(current page.center), inner sep=0pt] {
      \includegraphics[width=\paperwidth, height=\paperheight]{logos/background_figure.jpg}
    };
  \end{tikzpicture}
\end{frame}

% Resmi Baslik Sayfasi (kapak icin ozel alt bilgi)
{
  \setbeamertemplate{footline}{
    \begin{beamercolorbox}[wd=\paperwidth, ht=10pt, dp=1pt]{footline}
      \usebeamerfont{footline}
      \begin{tikzpicture}[remember picture,overlay]
        \node[anchor=south west, xshift=10mm, yshift=.4mm, align=left]
        at (current page.south west)
        {\tiny Bu şablon OTU Yazılım Mühendisliği tarafından hazırlanmış olup \href{https://github.com/ufukasia/Ostim-Teknik--niversitesi-Yaz-l-m-M-hendisli-i---Lutbeamer-egitim-i-eri-i-haz-rlama-}{buradan} güncel haline ulaşabilirsiniz.};
      \end{tikzpicture}
      \begin{tikzpicture}[remember picture,overlay]
        \node[anchor=south east, xshift=-5mm, yshift=0.5mm]
        at (current page.south east)
        {\insertframenumber \,/\, \inserttotalframenumber};
      \end{tikzpicture}
    \end{beamercolorbox}
  }
  \maketitle{}
}

% GENEL ICERIK HARITASI
\begin{frame}{İçindekiler}
  \scriptsize
  \begin{columns}[T]
    \begin{column}{0.48\textwidth}
      \tableofcontents[sections={1-4}]
    \end{column}
    \begin{column}{0.48\textwidth}
      \tableofcontents[sections={5-10}]
    \end{column}
  \end{columns}
\end{frame}

% ===================================================================
% ICERIK: HAFTA 7 - BEKLENEN DEGER, VARYANS, STANDART SAPMA
% ===================================================================

\section{Beklenen Değer Kavramı}
\subsection{Uzun Dönem Ortalama Fikri}

\begin{frame}{Deneysel Ortalama vs. Teorik Ortalama}
\end{frame}

\begin{frame}{Beklenen Değerin Yorumu}
\end{frame}

\subsection{Ayrık ve Sürekli Durumda Tanım}

\begin{frame}{Ayrık Rastgele Değişken için Beklenen Değer}
\end{frame}

\begin{frame}{Sürekli Rastgele Değişken için Beklenen Değer}
\end{frame}

\section{Fonksiyonların Beklenen Değeri}
\subsection{g(X) Kavramı}

\begin{frame}{Rastgele Değişken Üzerinden Yeni Rastgele Değişken}
\end{frame}

\begin{frame}{Dönüşüm ve Beklenen Değer}
\end{frame}

\subsection{Genel Beklenen Değer Teoremi}

\begin{frame}{Ayrık Durumda E[g(X)]}
\end{frame}

\begin{frame}{Sürekli Durumda E[g(X)]}
\end{frame}

\section{Birden Fazla Rastgele Değişken için Beklenen Değer}
\subsection{E[g(X, Y)] Kavramı}

\begin{frame}{Ortak Dağılım Üzerinden Beklenen Değer}
\end{frame}

\begin{frame}{Çok Boyutlu Beklenti Yorumu}
\end{frame}

\subsection{Marjinal Dağılımlar ile İlişki}

\begin{frame}{E(X) ve E(Y)'nin Marjinal Üzerinden Hesabı}
\end{frame}

\begin{frame}{Tek Değişkenli Beklentiye İndirgeme}
\end{frame}

\section{Varyans Kavramı}
\subsection{Yayılım Ölçüsü Olarak Varyans}

\begin{frame}{Ortalama Etrafında Dağılım}
\end{frame}

\begin{frame}{Sapmaların Karesi Fikri}
\end{frame}

\subsection{Varyansın Tanımı}

\begin{frame}{Var(X) Tanımı}
\end{frame}

\begin{frame}{$E[(X-\mu)^2]$ Yorumu}
\end{frame}

\section{Standart Sapma}
\subsection{Ölçekle Uyumlu Yayılım Ölçüsü}

\begin{frame}{Standart Sapmanın Tanımı}
\end{frame}

\begin{frame}{Varyans ile İlişkisi}
\end{frame}

\section{Beklenen Değer ve Varyansın Özellikleri}
\subsection{Lineerlik Özellikleri}

\begin{frame}{E(aX + b) Yapısı}
\end{frame}

\begin{frame}{Sabit ve Ölçekleme Etkisi}
\end{frame}

\subsection{Varyansın Dönüşüm Özellikleri}

\begin{frame}{Var(aX + b) Yapısı}
\end{frame}

\begin{frame}{Ölçeklemenin Etkisi}
\end{frame}

\section{Haftanın Kavramsal Özeti}
\subsection{Merkez ve Yayılım Perspektifi}

\begin{frame}{Beklenen Değer = Merkez}
\end{frame}

\begin{frame}{Varyans ve Standart Sapma = Yayılım}
\end{frame}

\subsection{İstatistiksel Modellemede Rolü}

\begin{frame}{Parametre Kavramına Geçiş}
\end{frame}

\begin{frame}{Tahmin ve Çıkarım İçin Zemin}
\end{frame}


\begin{frame}[allowframebreaks]
  \frametitle{Kaynaklar}
  \nocite{walpole_ch4}
  \printbibliography{}
\end{frame}

\appendix

\section{Yapay Zeka Asistanı}
\begin{frame}{Yapay Zeka Asistanınız: NotebookLM}
    \begin{columns}[T]
        \begin{column}{0.48\textwidth}
            \justifying{}
            \textbf{Ders Materyalleri ile Etkileşime Geçin!}

            Bu dersin tüm içeriği (sunumlar, makaleler ve notlar), Google'ın gelişmiş yapay zeka asistanı \textbf{NotebookLM} sistemine yüklenmiştir.

            \vspace{1em}
            \begin{itemize}
                \item \textbf{Anında Cevaplar:} Aklınıza takılanları sorun.
                \item \textbf{Özetler:} Karmaşık konuları basitçe öğrenin.
                \item \textbf{Audio Overview:} Dersi bir podcast gibi dinleyin.
            \end{itemize}

            \vspace{1.5em}
            \begin{center}
                {\setbeamercolor{button}{bg=red,fg=black}
                \href{https://notebooklm.google.com/notebook/1daaa713-b9d0-4bb3-bbe6-defbaf4e67d4}{\beamerbutton{\rule[-4ex]{0pt}{11ex}\Large \textbf{Hemen Deneyin}}}}
            \end{center}
        \end{column}

        \begin{column}{0.48\textwidth}
            \centering
            \includegraphics[width=\textwidth, height=0.6\textheight, keepaspectratio]{figures/notebooklm.jpg}
        \end{column}
    \end{columns}
\end{frame}

\begin{frame}[plain,noframenumbering]
    \begin{tikzpicture}[remember picture,overlay]
        \node[anchor=center, at=(current page.center), inner sep=0pt] {
            \includegraphics[width=\paperwidth, height=\paperheight]{figures/combined_qrs.png}
        };
        \node[anchor=south, at=(current page.south), yshift=0.5cm, fill=white!95, rounded corners, inner sep=10pt, text width=0.95\paperwidth, align=center] {
            \Large \textbf{\textcolor{blue}{Bizi Takip Edin!}} \\
            \vspace{0.3em}
            \small Ders tekrarlarıyla konuları pekiştirmek (YouTube), profesyonel proje ağlarına katılmak (LinkedIn) ve kampüs duyurularından anında haberdar olmak (Instagram) için topluluğumuza katılın.
        };
    \end{tikzpicture}
\end{frame}
\end{document}
